%% SCRAP: scratch/src/internal/PHYSICS_IMPLEMENTATION_OPTIONS
%% SOURCE: docs/working/scratch/src/internal/PHYSICS_IMPLEMENTATION_OPTIONS.adoc
%% STATUS: HISTORICAL
%% FITS: dev-guide/ch-physics
%% EDITORIAL: lifted — prose rewritten to press voice

\section{Physics Engine Implementation Options (Pre-Implementation Survey)}

This document records the three implementation paths considered for
extending the StarForth physics engine beyond its Phase 1 infrastructure.
It predates the current implementation and is HISTORICAL; the physics engine
has since been developed. Preserved as a design rationale record.

\subsection{State at Survey Time}

Phase 1 infrastructure was complete:

\begin{itemize}
  \item Execution hooks wired at both outer interpreter
        (\texttt{vm.c:533}) and colon word executor (\texttt{vm.c:456})
  \item Word metadata initialized with a 20-word seed table
        (\texttt{physics\_metadata\_apply\_seed})
  \item Temperature via exponential moving average from entropy
  \item Analytics heap allocated (10\,MiB default)
  \item Host snapshots functional (POSIX comprehensive, L4Re stub)
\end{itemize}

Missing: periodic capture mechanism, temperature decay model, feedback loop
integration.

\subsection{Option A: Zone 1 Discovery (1--2 hours)}

Wire the observable data stream end-to-end to verify metrics flow from word
execution to analytics heap. Steps: (1) call \texttt{physics\_host\_snapshot()}
at 10\,ms intervals in the REPL loop; (2) add a \texttt{PHYSICS-HEAP-DUMP}
FORTH word to inspect heap contents; (3) add a \texttt{PHYSICS-TEMPS} word
to display per-word temperature distribution; (4) write an end-to-end test
confirming metrics flow.

\subsection{Option B: Metrics Infrastructure (2--3 days)}

Build comprehensive metrics collection to feed scheduling and memory decisions.
Steps: (1) per-word call-stack depth and recursion tracking; (2) latency
histogram with percentiles (P50, P99, P99.9); (3) call-graph discovery with
edge frequencies; (4) composite summary words (\texttt{PHYSICS-SUMMARY}).

\subsection{Option C: Design Questions Resolution (1 day)}

Resolve six architectural decisions before committing to a control system:

\begin{description}
  \item[Scheduling model.] Option chosen: hybrid (advisory hints building
        toward batch enforcement).
  \item[Memory placement.] Option chosen: logical tiers (advisory labels,
        not physical movement).
  \item[GC vs.\ spilling.] Option chosen: governance-driven policy with
        per-word \texttt{can\_forget} and \texttt{can\_spill} flags.
  \item[L4Re specifics.] Decision: defer L4Re optimization to Phase 3;
        complete POSIX path first.
  \item[Governance integration.] Decision: load governance config at
        startup; runtime adjustments stay within approved ranges.
  \item[Measurement overhead.] Estimated additional cost 80--100\,ns per
        decision; acceptable for interactive and high-performance targets.
\end{description}

\subsection{Recommended Sequence}

A $\to$ C $\to$ B $\to$ Phase 2 implementation. Option A validates
infrastructure quickly; Option C guides what metrics Option B should
prioritize; Option B then provides the data required for Phase 2 weak
control.
